% !TEX root = ../main.tex
ソフトウェア開発では、私たちは日常的に「変更しても意味が変わっていないか」を気にしています。
関数の分割、APIレスポンスの形の変更、DBスキーマの移行、古いデータから新しい形式への変換、複数のサービスから得た情報の結合などが、その例です。
これらは一見ばらばらの作業に見えますが、「ある形のものを別の形へ変換し、その変換が期待した性質を保つか」を確認する点で共通しています。

圏論は、このような変換と合成を扱うための言語です。
対象、射、合成、恒等射、同型、関手、自然変換といった言葉は、最初は抽象的に見えます。
しかし、ソフトウェアの視点では、対象を型やデータ構造、射を関数やマイグレーション、合成をパイプライン、同型を情報を失わない相互変換、関手を構造を保った一括変換として読めます。

一方、Lean 4 は「この変換は本当に正しいのか」を機械的に確認するための道具です。
通常のテストは具体例を確認し、型は多くの不正な組み合わせを排除します。
しかし、「すべての入力についてこの性質が成り立つ」「このマイグレーションは戻せる」「このリファクタリングは振る舞いを変えない」といった主張は、テストだけでは原理的に確認しきれません。
Lean では、そのような主張を命題として書き、証明として検査できます。

この部では、本書全体の地図を示します。
圏論は研究対象としてではなく、設計、仕様、変更、同値性を整理するための道具として扱います。
Lean も本番コードを置き換えるものではなく、仕様の核心部分を小さなモデルとして表し、変更前後の意味保存を確認するための実験室として扱います。

この部を読み終えるころには、本書で扱う中心的な問いを把握できます。
その問いは、「この処理は何と何を変換しているのか」「どの性質を保つべきなのか」「それはテストで確認すべきか、型で表すべきか、Leanで証明すべきか」の三つです。

\section*{この部で扱うこと}

\begin{itemize}
\item ソフトウェアにおける正しさの種類
\item テスト、型、仕様、証明の役割分担
\item 圏論を合成と変換の言語として見る視点
\item Lean 4 を仕様確認の道具として使う視点
\item 本書で扱う実務サンプルの全体像
\end{itemize}

\section*{この部で扱わないこと}

\begin{itemize}
\item 圏論の厳密な公理化
\item Lean 4 の詳細な構文説明
\item Mathlib の高度な使い方
\item 実システム全体の完全検証
\end{itemize}

\section*{次の部への接続}

次の部では、Lean 4 の最小限の読み書きを学びます。
最初から難しい数学を形式化するのではなく、型、値、関数、命題、証明という基本要素を使って、「仕様をコードの形で書く」とはどういうことかを確認します。
